\section*{الفصل \ref{ch:b2:chromatography} --- الكروماتوغرافيا}

\begin{solution}{exo:b2:chromatography:1}
$k = (5.0 - 1.0)/1.0 = 4.0$؛ وكسر الزمن في \omterm{def:b2:chromatography:principle}{الطور المتحرك} $1/(1 + k) = 0.20$.
\end{solution}

\begin{solution}{exo:b2:chromatography:2}
$N = 16(6.0/0.30)^2 = 16 \times 400 = 6400$.
\end{solution}

\begin{solution}{exo:b2:chromatography:3}
$H = \qty{250}{mm}/10\,000 = \qty{0.025}{mm} = \qty{25}{\micro m}$.
\end{solution}

\begin{solution}{exo:b2:chromatography:4}
الإيثانول في الدم: GC (متطاير). البروتين: HPLC. الكافيين في مشروب: HPLC (في الماء، وليس متطايرًا
بما يكفي دون تحضير). البنزين في وقود السيارات: GC. السكريات: HPLC (تتفكك قبل أن تغلي).
\end{solution}

\begin{solution}{exo:b2:chromatography:5}
$R_s = 2(4.6 - 4.0)/(0.40 + 0.50) = 1.2/0.90 \approx 1.3$: ليستا مفصولتين تمامًا حتى خط الأساس (دون 1.5).
\end{solution}

\begin{solution}{exo:b2:chromatography:6}
$\sqrt N = 4R_s\,\dfrac{\alpha}{\alpha - 1}\,\dfrac{1 + k_2}{k_2} = 4 \times 1.5 \times 11 \times 1.25 =
82.5$، إذن $N \approx 6.8 \times 10^3$.
\end{solution}

\begin{solution}{exo:b2:chromatography:7}
$u_{\mathrm{opt}} = \sqrt{8/2} = \qty{2}{mm/s}$؛ $H_{\min} = 5 + 2\sqrt{16} = \qty{13}{\micro m}$.
\end{solution}

\begin{solution}{exo:b2:chromatography:8}
الفينول (\ce{OH} قطبية، وروابط هيدروجينية مع الماء) أولًا، ثم البنزين، ثم التولوين (بمجموعة
\ce{CH3} إضافية، فهو أقل قطبية). وزيادة الميثانول تجعل \omterm{def:b2:chromatography:principle}{الطور المتحرك} أقل قطبية: فتخرج الثلاثة أبكر، بالترتيب
نفسه.
\end{solution}

\begin{solution}{exo:b2:chromatography:9}
$F = (860/400)/(20.0/10.0) = 2.15/2.00 = 1.075$. العينة: $c_X = (645/410) \times 10.0/1.075 \approx
\qty{14.6}{mg/L}$.
\end{solution}

\begin{solution}{exo:b2:chromatography:10}
$R_s \propto \sqrt N \propto \sqrt L$: يجب ضرب $L$ في $(1.5/1.06)^2 \approx 2.0$، مما يضاعف
زمن التحليل والضغط. والروافع الأخرى: الانتقائية (\omterm{def:b2:chromatography:principle}{الطور المتحرك}، \omterm{def:b2:chromatography:principle}{والطور الثابت}، ودرجة
الحرارة)، والجسيمات الأصغر، ومعدل التدفق الأمثل.
\end{solution}

\begin{solution}{exo:b2:chromatography:11}
تفصل بين القمتين مسافة $4\sigma R_s$، فتبعد نقطة المنتصف $2\sigma R_s$ عن كل منهما. عند $R_s = 1$: تساهم كل قمة
بالمقدار $\mathrm e^{-(2)^2/2} = \mathrm e^{-2} = 0.135$؛ فيكون الوادي عند $0.27$ من ارتفاع القمة.
وعند $R_s = 1.5$: $2\mathrm e^{-(3)^2/2} = 2\mathrm e^{-4.5} \approx 0.022$، أي نحو 2~\%: عودة إلى خط الأساس
من الناحية العملية.
\end{solution}

\begin{solution}{exo:b2:chromatography:12}
$k_2/(1+k_2)$: 0.67، 0.83، 0.91؛ والزمن بوحدة $t_M$: 3، 6، 11. فالانتقال من 2 إلى 5 يربح 25~\% في
\omterm{def:b2:chromatography:separation}{الميز} مقابل ضعف الزمن؛ ومن 5 إلى 10، يربح 9~\% مقابل نحو ضعف الزمن مرة أخرى. فقيمة $k$ بين نحو 2 و5 هي
الحل الوسط المعتاد.
\end{solution}

\begin{solution}{pb:b2:chromatography:1}
\textbf{1.} الثابت: حبيبات سيليكا تحمل سلاسل ألكيلية طويلة، غير قطبية. والمتحرك: ماء مع
الميثانول، قطبي.
\textbf{2.} إنها شديدة القطبية وتبقى في \omterm{def:b2:chromatography:principle}{الطور المتحرك}: $k \approx 0$.
\textbf{3.} للكافيين مجموعة ميثيل واحدة أكثر من الثيوفيلين، وللثيوفيلين رابطة \ce{N-H}
تكوّن روابط هيدروجينية مع الماء: فالثيوفيلين أكثر قطبية ويخرج أولًا.
\textbf{4.} فهو قريب البنية (استجابة وسلوك متشابهان)، وغائب عن المشروب، ويخرج
قرب الكافيين مع بقائه مفصولًا عنه.
\textbf{5.} سينقص كلاهما: \omterm{def:b2:chromatography:principle}{فالطور المتحرك} الأقل قطبية ينافس \omterm{def:b2:chromatography:principle}{الطور الثابت} على نحو أفضل.
\textbf{6.} الزمن الذي يستغرقه \omterm{def:b2:chromatography:principle}{الطور المتحرك} لعبور العمود، وهو الزمن الذي يقضيه المركب متحركًا.
\textbf{7.} $k_{\mathrm{theo}} = (3.0 - 1.0)/1.0 = 2.0$؛ $k_{\mathrm{caf}} = 2.4$.
\textbf{8.} $\alpha = 2.4/2.0 = 1.2$.
\textbf{9.} الكافيين $N = 16(3.4/0.20)^2 = 4624 \approx 4.6 \times 10^3$؛ والثيوفيلين $16(3.0/0.18)^2
\approx 4.4 \times 10^3$.
\textbf{10.} $H = \qty{150}{mm}/4624 \approx \qty{0.032}{mm} = \qty{32}{\micro m}$.
\textbf{11.} $R_s = 2(3.4 - 3.0)/(0.18 + 0.20) = 0.80/0.38 \approx 2.1$.
\textbf{12.} $\frac{\sqrt{4624}}{4} \times \frac{0.2}{1.2} \times \frac{2.4}{3.4} = 17 \times 0.167 \times
0.706 \approx 2.0$؛ والفرق الصغير يأتي من افتراض تساوي العرضين في المعادلة.
\textbf{13.} نعم: $R_s > 1.5$.
\textbf{14.} يُنصَّف $N$: $R_s \approx 2.1/\sqrt2 \approx 1.5$، عند حد الفصل حتى خط الأساس بالضبط.
\textbf{15.} يُنصَّف: الكافيين عند \qty{1.7}{min}.
\textbf{16.} $R_s \approx 2.1\sqrt2 \approx 3.0$، مقابل ضعف الزمن وضعف الضغط: وهذا هدر هنا.
\textbf{17.} تركيب \omterm{def:b2:chromatography:principle}{الطور المتحرك} (كسر الميثانول، وpH، ومذيب عضوي آخر) أو
\omterm{def:b2:chromatography:principle}{الطور الثابت}.
\textbf{18.} قليلًا: فقيمة $k_2/(1 + k_2) = 0.71$ أصلًا؛ وعند $k_2 = 5$ تصير 0.83 (+18~\%) مقابل
زمن تحليل مضروب في $6/3.4 \approx 1.8$.
\textbf{19.} $F = (1500/1000)/(50.0/40.0) = 1.50/1.25 = 1.20$.
\textbf{20.} القمتان كلتاهما من الحقنة نفسها: فيُختزل الحجم في نسبة المساحتين.
\textbf{21.} $c = (970/1010) \times 40.0/1.20 \approx \qty{32.0}{mg/L}$.
\textbf{22.} التخفيف عشر مرات: \qty{320}{mg/L}.
\textbf{23.} ستكون $A_{IS}$ أكبر مما ينبغي، ونتيجة الكافيين أدنى مما ينبغي.
\textbf{24.} سلسلة من معايير الكافيين تُحقن في الظروف نفسها بالضبط التي تُحقن فيها العينة، 
وخط معايرة للمساحة بدلالة التركيز، وحجوم حقن قابلة للتكرار.
\textbf{25.} $\qty{320}{mg/L} \times \qty{0.250}{L}$: $\boldsymbol{\approx \qty{80}{mg}}$ من الكافيين لكل علبة (معطيات
التمرين).
\end{solution}
